\documentclass[12pt]{article}
\usepackage[utf8]{inputenc}
\usepackage[T1]{fontenc}
\usepackage[margin=1in]{geometry}
\usepackage{setspace}
\usepackage{xcolor}
\usepackage[colorlinks=true, urlcolor=blue!60!black, linkcolor=blue!60!black]{hyperref}
\setlength{\parindent}{2em}
\setlength{\parskip}{1.25em}
\renewcommand{\baselinestretch}{1.0}
\begin{document}
\begin{center}
{\LARGE \textit{Roe v. Wade}, 410 U.S. 113 (1973)\par}
\vspace{0.8em}
{\large Case Brief --- Four-Page Edition\par}
\end{center}
\vspace{1em}

\section{Memory Jogger}

The Fourteenth Amendment's Due Process Clause protects a fundamental right of personal privacy broad enough to encompass a woman's decision to terminate a pregnancy, governed by a trimester framework balancing maternal health and fetal viability; the framework was later modified in \textit{Casey} (1992), and \textit{Roe} was overruled in \textit{Dobbs} (2022) (\href{https://supreme.justia.com/cases/federal/us/410/113/}{410 U.S. 113}).

\section{Detailed Case Facts}

In March 1970, Norma McCorvey, suing as ``Jane Roe,'' was an unmarried, pregnant Dallas County resident who wanted a safe abortion by a licensed physician. Texas statutes dating to 1854 made procuring an abortion a felony unless done on medical advice to save the mother's life. Her life was not threatened, and she could not afford to travel elsewhere.

Represented by Linda Coffee and Sarah Weddington, Roe brought a class action against Dallas District Attorney Henry Wade. Dr.\ James Hallford, who faced two pending abortion prosecutions, intervened, and John and Mary Doe, a childless married couple fearing contraceptive failure, filed a companion complaint. Roe gave birth before the case was resolved, and Texas argued mootness. The Court heard the case with \textit{Doe v.\ Bolton}, which challenged Georgia's statute.

\section{Procedural History}

A three-judge district court held the statutes void for vagueness and overbroad under the Ninth Amendment, dismissed the Does, and denied an injunction on abstention grounds. Both sides appealed directly. The case was argued in December 1971, reargued in October 1972 after Justices Powell and Rehnquist joined the Court, and decided on January 22, 1973, affirmed in part and reversed in part.

\section{Judicial Votes}

\textbf{7--2} to invalidate the Texas statutes.
\begin{itemize}
\item \textbf{Majority:} Justice Blackmun, joined by Chief Justice Burger and Justices Douglas, Brennan, Stewart, Marshall, and Powell.
\item \textbf{Concurring:} Chief Justice Burger; Justice Douglas; Justice Stewart.
\item \textbf{Dissenting:} Justice White, joined by Justice Rehnquist; Justice Rehnquist.
\end{itemize}

\section{Holding}

Pregnancy is ``capable of repetition, yet evading review,'' so Roe's claim was not moot; Hallford's claim was dismissed under \textit{Younger}, and the Does lacked standing. Due process protects a fundamental privacy right that includes a qualified abortion decision, subject to strict scrutiny, and a fetus is not a constitutional ``person.'' Under the trimester framework, the first-trimester decision belongs to the woman and her physician; afterward the State may regulate to protect maternal health; and after viability it may proscribe abortion except where necessary to preserve the mother's life or health. A State may bar abortions by non-physicians. The historical surveys and examples of permissible regulations were dicta.

\section{Analysis of Opinions}

\textbf{Majority (Blackmun).} After resolving justiciability, Blackmun surveyed ancient attitudes, the Hippocratic Oath, the common law's quickening distinction, nineteenth-century statutes, and professional associations, concluding that restrictive bans were of recent vintage. Modern medicine made early abortion at least as safe as childbirth, weakening the health rationale, though the State had a legitimate interest in prenatal life. Drawing on \textit{Meyer}, \textit{Pierce}, \textit{Skinner}, \textit{Griswold}, \textit{Loving}, and \textit{Eisenstadt}, he located the privacy right in Fourteenth Amendment liberty, noting the medical, psychological, and financial burdens of unwanted pregnancy. Constitutional uses of ``person'' apply postnatally, and Texas did not treat abortion as murder. Declining to decide when life begins, the Court balanced interests by trimester, and the Texas statutes fell as a unit.

\textbf{Burger, Douglas, and Stewart (concurring).} Burger stressed that the rulings did not establish abortion on demand. Douglas read the Ninth Amendment's retained rights as informing Fourteenth Amendment liberty. Stewart conceded that \textit{Griswold} was a substantive due process decision.

\textbf{White (dissenting).} Nothing in the Constitution's language or history supports the right; the Court's judgment was ``an exercise of raw judicial power'' over a moral question for the people.

\textbf{Rehnquist (dissenting).} The record did not show Roe was in her first trimester; a surgical operation is not ``private''; rational-basis review should apply, as under \textit{Lee Optical}; the approach echoes \textit{Lochner}; and at least 36 states or territories restricted abortion in 1868.

\section{Examples: Future Applications}

\begin{enumerate}
\item \textbf{Same-Side (Invalid): Hospital-Only Rule.} A law requires all abortions, even at eight weeks, to occur in hospitals. The first-trimester decision is free of state regulation beyond requiring a licensed physician.
\item \textbf{Same-Side (Invalid): No Health Exception.} A post-viability ban lacks an exception for serious health threats, though \textit{Roe} requires one for the mother's life or health.
\item \textbf{Opposite-Side (Upheld): Non-Physicians.} A state prosecutes an unlicensed practitioner; \textit{Roe} allows proscribing non-physician abortions at any stage.
\item \textbf{Opposite-Side (Upheld): Second-Trimester Safety Rules.} Clinic equipment, transfer, and staffing rules at 16 weeks are reasonably related to maternal health.
\item \textbf{Fence-Sitter: Waiting Periods.} A 24-hour waiting period and state script at nine weeks arguably intrude on the physician's first-trimester judgment yet resemble ordinary informed consent; \textit{Akron} struck such rules, and \textit{Casey} later upheld one.
\end{enumerate}

\section{Critique}

\textbf{Originalist critiques.} Echoing Rehnquist, critics note that the Constitution mentions neither privacy nor abortion and that 36 states and territories restricted abortion in 1868. In \textit{Dobbs}, Justice Alito argued that \textit{Roe}'s history was flawed, that its trimester scheme read like a statute, and that it short-circuited democratic self-government.

\textbf{Progressive and feminist critiques.} John Hart Ely, though he favored legal abortion, wrote that \textit{Roe} was bad constitutional law. Ruth Bader Ginsburg argued that the Court moved too far too fast, rather than striking only Texas's extreme statute, and grounded the right in medicalized privacy rather than equal protection and women's equal citizenship. Feminist scholars note that the opinion centered the attending physician's judgment rather than women's autonomy.

\section{Key Quotations}

\begin{quote}
``This right of privacy, whether it be founded in the Fourteenth Amendment's concept of personal liberty and restrictions upon state action, as we feel it is, or, as the District Court determined, in the Ninth Amendment's reservation of rights to the people, is broad enough to encompass a woman's decision whether or not to terminate her pregnancy.''
\end{quote}
\noindent --- Justice Blackmun, Opinion of the Court, 410 U.S. at 153.

\begin{quote}
``All this, together with our observation, \textit{supra}, that, throughout the major portion of the 19th century, prevailing legal abortion practices were far freer than they are today, persuades us that the word `person,' as used in the Fourteenth Amendment, does not include the unborn.''
\end{quote}
\noindent --- Justice Blackmun, Opinion of the Court, 410 U.S. at 158.

\begin{quote}
``We need not resolve the difficult question of when life begins. When those trained in the respective disciplines of medicine, philosophy, and theology are unable to arrive at any consensus, the judiciary, at this point in the development of man's knowledge, is not in a position to speculate as to the answer.''
\end{quote}
\noindent --- Justice Blackmun, Opinion of the Court, 410 U.S. at 159.

\end{document}
